\section{Experimental methodology and results}
\label{sec:experiments}

\subsection{Evidence hierarchy and environment}

The evaluation follows the evidence hierarchy in Table~\ref{tab:evidence}. C2 is the core correctness study. A1 and P1 bound model-side repair cost and the prepared-state software control path. E3 probes a downstream automated outcome. C-E2 and C-E1 characterize candidate generation in a separate synchronous harness, while A2 is exploratory. These campaigns have different states and timing boundaries; their absolute times are neither pooled nor subtracted across campaigns.

\begin{table}[t]
\centering
\small
\caption{Evidence hierarchy and measurement boundary.}
\label{tab:evidence}
\begin{tabularx}{\linewidth}{@{}lYY@{}}
\toprule
Study & Role & Principal limitation \\
\midrule
C2 v3 & Direct crop and recovery correctness & One snapshot/backend; no clean re-prefill \\
A1/P1 & Fixed-protocol cost and software path & No device or acoustic stopping \\
E3 & Automated downstream probe & No human reference standard \\
C-E2 & Token-consistent candidate characterization & Synchronous oracle \\
C-E1 & Whole implementation-path audit & Paths not token-equivalent \\
A2 & Exploratory history naturalization & Generation trajectories confounded \\
\bottomrule
\end{tabularx}
\end{table}

Data were derived from MultiWOZ 2.1 \citep{eric2020MultiWOZ21}. The principal model was Qwen2-7B-Instruct \citep{yang2024Qwen2}; TEN Turn Detection supplied turn confidence \citep{tenFramework2025TurnDetection}; Mistral-7B-Instruct-v0.3 supplied the E3 automated judgments \citep{jiang2023Mistral7B,mistralAI2024Mistral7BInstructV03}; a six-sentence CosyVoice2-0.5B profile parameterized deterministic Mock-TTS durations \citep{du2024CosyVoice2}; and Qwen3-0.6B was used by the exploratory rewriter \citep{yang2025Qwen3}. The implementation reused Whisper-based ASR \citep{radford2023RobustSpeechRecognition} and stream2sentence version 1.0.0 for output segmentation \citep{beigel2026Stream2Sentence}.

The frozen runtime used Python 3.10.18, PyTorch 2.8.0+cu128, Transformers 4.57.1, CUDA 12.8, and cuDNN 9.10.2 on RTX 3090 GPUs with 24~GB memory. C2 used BF16 and SDPA. Confirmatory C-E1/C-E2 used two GPUs; other campaigns followed their own manifests.

\subsection{C2 exact correctness}

C2 v1 and v2 retain their rejected verdicts. The v2 structural checks passed, but its numerical gate passed only 42/45 comparisons and its clean-reprefill control did not match the production forward topology. It could therefore establish neither crop attribution nor clean-reprefill equivalence. C2 v3 posed the narrower, identifiable same-snapshot slicing and matched-recovery questions.

\begin{table}[t]
\centering
\small
\caption{C2 v3 exact-gate coverage.}
\label{tab:c2}
\begin{tabularx}{\linewidth}{@{}Yr@{}}
\toprule
Check & Result \\
\midrule
Cases / crop events & 24/24; 27/27 \\
Production assistant append & 308 tokens \\
K/V layers compared & 28 \\
Recovery steps & 60/60 \\
Wrong-length controls detected & 27/27 \\
No-op crops & 3/3 \\
\bottomrule
\end{tabularx}
\end{table}

For all 27 crop events, the retained pre-crop K/V prefix, production post-crop K/V, and slicing-oracle K/V agreed across 28 layers in shape, dtype, device, hash, and runtime \texttt{torch.equal}. The independently derived keep length, mask, global ledger, sequence length, and K/V length were exact. Across 60 recovery steps, matched arms remained exact in K/V, logits, mask, ledger, retained-prefix hash, and role/end/content state. All wrong-length controls were detected.

These results support direct crop integrity and within-run matched-arm recovery exactness for the tested configuration. They do not revise the rejected v1/v2 verdicts and do not establish clean-reprefill numerical equivalence, continuation equivalence, cross-process or cross-device determinism, or end-to-end production correctness.

\subsection{Fixed-protocol repair cost and software path}

A1 compared synchronized joint crop-plus-role repair with complete re-prefill over 256--8192-token contexts. Each length used five warm-ups and 50 recorded repetitions, fixed operation order, and removal of a fixed 32-token suffix. Table~\ref{tab:a1} and Figure~\ref{fig:a1} show that joint-repair medians ranged from 31.054 to 48.315~ms. The median re-prefill-to-repair ratio increased from $2.254\times$ at 256 tokens to $40.620\times$ at 8192 tokens. This is a latency microbenchmark, not a state-equivalence test.

\begin{table*}[t]
\centering
\small
\caption{A1 synchronized repair microbenchmark; 50 repetitions per length.}
\label{tab:a1}
\begin{tabular*}{\textwidth}{@{\extracolsep{\fill}}lrrrrrr@{}}
\toprule
Context length & 256 & 512 & 1024 & 2048 & 4096 & 8192 \\
\midrule
Joint median (ms) & 31.616 & 31.852 & 31.054 & 31.519 & 36.903 & 48.315 \\
Joint IQR (ms) & 2.356 & 2.162 & 3.099 & 1.197 & 0.635 & 0.928 \\
Re-prefill/joint & $2.254\times$ & $4.124\times$ & $7.707\times$ & $15.020\times$ & $25.453\times$ & $40.620\times$ \\
\bottomrule
\end{tabular*}
\end{table*}

\begin{figure}[t]
\centering
\includegraphics[width=0.82\linewidth,height=0.39\textheight,keepaspectratio]{Figure_5_cost_microbenchmark.pdf}
\caption{A1 joint repair and re-prefill microbenchmark. Error bars show interquartile ranges across 50 repetitions.}
\label{fig:a1}
\end{figure}

P1 crossed three context lengths (512, 2048, and 8192 tokens) with three cursor proportions (0.25, 0.50, and 0.75) and 20 repetitions per cell, for 180 prepared-state headless records. All 180 requests and acknowledgements reached the intended software cursor. A zero leaked-sample count means only that the software counter did not advance after acknowledgement. Cell medians were 0.055--0.062~ms for stop acknowledgement, 0.167--0.176~ms for post-acknowledgement GPU synchronization, 0.47--0.50~ms for timeline lookup, 2.44--2.53~ms from stop to crop completion, and 78.6--80.8~ms from stop to role recovery. The latter intervals are nested. With $n=20$ per cell, empirical P95 values are descriptive, not service-level objectives. Neither A1 nor P1 includes device/acoustic stopping, online TTS cancellation, or a complete concurrent barge-in path.

\subsection{Downstream fixed-trajectory probe}

E3 contained 100 dialogues, 400 paired dialogue--injection-position units, 800 condition records, and 1600 judge records. Each assistant trajectory was generated once; software-cursor and generation-retention conditions shared that trajectory, its fragment timeline, and four injection positions (0.25, 0.50, 0.75, and a fragment boundary). Fragment and character-ratio proxy targets were separately eligible when their respective target strings were non-empty. The fragment analysis contained 297 eligible labels from 96 dialogues; the proxy analysis contained 380 labels from all 100 dialogues. Eligibility and uncertainty are therefore target-specific rather than based on 400 independent observations.

The primary estimand assigned equal weight to each eligible dialogue--position label and computed generation minus software-cursor retention. Confidence intervals used 10,000 paired dialogue-cluster bootstrap repetitions. A dialogue-weighted sensitivity first averaged within dialogue. Target-specific exact-key deduplication combined dialogue and trajectory identities, both history keys, and an exact target hash; the fragment key also included the retained token endpoint. It is not semantic clustering.

The lexical detector matched numbers, selected capitalized terms, and non-stopword content terms. The \texttt{specific-reference-v3} Mistral judge received the target and two delimiter-joined replies without condition identity, and greedily produced a strict YES/NO decision. Neither detector is a human reference standard. The intervals quantify dialogue-sampling uncertainty conditional on the frozen trajectories, targets, detectors, prompts, and 40-token response cap; they do not include detector error or model/prompt variation.

\begin{table}[t]
\centering
\small
\caption{E3 label-weighted information-reproduction rates. The effect is generation minus software-cursor retention.}
\label{tab:e3}
\begin{tabularx}{\linewidth}{@{}llccY@{}}
\toprule
Target & Detector & Cursor rate & Generation rate & Effect [95\% cluster CI] (pp) \\
\midrule
Fragment & Lexical & 67.00\% (199/297) & 63.64\% (189/297) & $-3.37$ [$-10.49$, 3.40] \\
Fragment & Judge & 42.76\% (127/297) & 40.74\% (121/297) & $-2.02$ [$-10.70$, 6.13] \\
Proxy & Lexical & 75.26\% (286/380) & 73.68\% (280/380) & $-1.58$ [$-6.08$, 2.67] \\
Proxy & Judge & 43.95\% (167/380) & 41.32\% (157/380) & $-2.63$ [$-8.57$, 2.90] \\
\bottomrule
\end{tabularx}
\end{table}

Figure~\ref{fig:e3} includes the main and sensitivity intervals. Exact-key deduplication reduced 297 fragment labels to 169 groups and 380 proxy labels to 379 groups. Dialogue weighting yielded effects from $-3.21$ to $-1.30$ percentage points; exact-key effects ranged from $-2.96$ to 0.00 percentage points. All corresponding intervals crossed zero. The four primary estimates therefore establish neither advantage nor harm, and they do not establish equivalence, non-inferiority, absence of a difference, or human-semantic effects.

\begin{figure}[t]
\centering
\includegraphics[width=\linewidth,height=0.34\textheight,keepaspectratio]{Figure_6_e3_intervals.pdf}
\caption{E3 generation-minus-software-cursor effects. Points and intervals show the label-weighted analysis and two sensitivity analyses.}
\label{fig:e3}
\end{figure}

\subsection{Candidate-generation characterization}

Confirmatory C-E1/C-E2 used 100 unique holdout utterances, five independently initialized process sessions, and ten conditions. Each condition contained 500 session--utterance observations, but the content sampling unit remained 100 utterances. Formal inference used 10,000 crossed/product bootstrap repetitions by independently resampling sessions and utterances and evaluating their Cartesian-product multiplicities.

The harness delivered pre-segmented text synchronously. It excluded real audio, ASR wall-clock time, online turn-detector computation, TTS, playback, and a sound card. \texttt{first\_token\_ready} followed candidate-token selection but preceded cache-update completion and generator yield, so it is reported as candidate selection/readiness. \texttt{endpoint\_accept} was a post-candidate synchronous oracle, not a natural endpoint detector. The effective-TTFT quantity is therefore an optimistic oracle lower bound.

C-E2 compared the token-consistent B@0.92 and never-trigger paths. Candidate selection/readiness remained between 62.0 and 62.4~ms across all nine B-path points. At B@0.92, candidate availability at oracle acceptance was 67\% (335/500), with crossed 95\% CI [58\%, 76\%], and the pooled discarded-token ratio was 2.85\% [1.12\%, 4.73\%]. The latter is a ratio of token counts, not FLOPs, GPU time, energy, or bandwidth. Relative to never trigger, the readiness difference (never minus B@0.92) was $-0.03$~ms [$-0.64$, 0.61], while the oracle-latency lower-bound difference was $+20.80$~ms [17.85, 23.65]. Across the nine discrete points, lower thresholds increased candidate availability and the pooled discarded-token ratio, while arrival-to-candidate readiness remained between 62.0 and 62.4~ms. These findings characterize generation before synchronous oracle acceptance; they do not show readiness before real end of speech.

C-E1 compared full-string tokenization/full prefill with segment-wise tokenization/incremental forward. The readiness means were 27.70 and 62.38~ms, respectively, giving an A-minus-B difference of $-34.69$~ms [$-35.44$, $-33.95$]. Oracle lower-bound means were 27.70 and 10.26~ms, for $+17.44$~ms [14.41, 20.32]. Complete token outputs matched in only 280/500 pairs and diverged for at least one session in 44/100 utterances. The contrast therefore combines tokenization, forward topology/shape, role boundaries, kernels, and Python scheduling. It cannot identify an incremental-prefill effect, and selecting only output-matched pairs would be post-treatment selection.

Finally, A2 generated 100 records per naturalization implementation. Mean automated coherence scores were 3.76 for naive retention, 3.62 for lightweight rewriting, and 3.29 for interruption marking. Rewriting took 639~ms on average (median 670~ms; interpolated P90 approximately 935~ms; maximum 1165~ms). Because the conditions regenerated both initial and subsequent responses and shared compatible retained text in only 33/100 dialogues, these values do not estimate a strategy effect. They are exploratory descriptions, not negative or null results.
